%============================================================
\chapter{Unmanned Operation: The Limit With the Operator Removed}
\label{part:unmanned}
%============================================================

\section{The question}

Chapter~\ref{ch:solo} treated a solo business. This chapter asks about its limit: whether a
business with no operator at all can stand. Since it is obvious that launching requires
people, the question of unmanned operation is confined to steady running, after which the
launch cost itself is examined to see how far it can be reduced.


\begin{center}
\begin{tabularx}{\textwidth}{@{}llX@{}}
\toprule
\multicolumn{3}{@{}l}{\textbf{Theoretical quantities in this chapter}}\\
\midrule
$\Omega,\ \hat\Omega$ & Chapter~\ref{sec:info} & the state space, and the subset envisaged at design time \\
$\lambda_{\rm novel}$ & this chapter & the arrival rate of $\Omega\setminus\hat\Omega$ \\
$\mathcal{G}$ & Chapter~\ref{sec:info} & automation demands a stronger condition than this \\
$\kappa_C>0$ & Eq.~\eqref{eq:kappa} & credit issued by the operator as a free period \\
\bottomrule
\end{tabularx}
\captionof{table}{Theoretical quantities in this chapter.}
\label{tab:unmanned-quantities}
\end{center}

\section{Formalizing ``unmanned''}

\subsection{Automation is a stronger condition than measurability}

Chapter~\ref{sec:info} made $\mathcal{G}$-measurability the condition for a contract to be
implementable. Automation is stronger.

\begin{definition}[Automation]
Writing the response down \emph{in advance} as a measurable function
$a:\Omega\to\mathcal{A}$ on the state space.
$\mathcal{A}$ is the set of available responses (a different quantity from fixed assets $A$,
distinguished by typeface).
\end{definition}

Proposition~\ref{prop:implementable} required $\pi$ to be $\mathcal{G}$-measurable.
Automation adds the requirement that the response be written down \emph{in advance}. A
contract need only be verifiable after the fact; automation does not have that luxury. An
algorithm is nothing other than a measurable function specified in advance. Thus
\begin{equation}
\text{implementable}\;:\;\pi\ \text{is}\ \mathcal{G}\text{-measurable}
\qquad\subsetneq\qquad
\text{automatable}\;:\;a\ \text{specified in advance on}\ \hat\Omega
\end{equation}

The difficulty is that $\Omega$ is not closed. Writing $\hat\Omega\subset\Omega$ for the
state space envisaged at design time, $a$ is undefined when
$\omega\in\Omega\setminus\hat\Omega$ arrives. Legal change, a discovered vulnerability, a
payment provider's change of terms, and unforeseen use are such cases.

\begin{proposition}[Unmanned operation is a function of duration]
\label{prop:unmanned-period}
Writing $\lambda_{\rm novel}$ for the arrival rate of unforeseen states, the period over
which operation can be unmanned is about $1/\lambda_{\rm novel}$. Fully unmanned operation
holds for a finite period but not indefinitely.
\end{proposition}

\begin{remark}[An upper bound on $\lambda_{\rm novel}$ can be measured]
\label{rem:lambda-upper}
In the ``inside'' partition of Chapter~\ref{ch:protocol}, \emph{responses} to unforeseen
states leave a record. Modification runs at 0.048 and cessation at 0.030 a year, and
arrivals requiring no response cannot be counted, so these are a lower bound on
$\lambda_{\rm novel}$ (Section~\ref{sec:protocol-byproducts},
Remark~\ref{rem:lambda-is-response}). The $1/\lambda_{\rm novel}$ of this proposition is
therefore shorter than 20 years.

The distribution is bimodal. Those that are touched are touched with a median of 42 days;
those that are not go untouched for a median of 4.6 years.
\textbf{A mean $1/\lambda$ conceals this bimodality.}
\end{remark}

\subsection{What is required is not labour but availability}

Responding to $\Omega\setminus\hat\Omega$ is not steady labour. Normally nothing is done, and
action occurs only on arrival. The structure is the same as family 2-5 (options and
warranties) of Part~\ref{part:catalog}.

\begin{equation}
\text{the operator of an unmanned business is short an option on }\Omega
\end{equation}

Delivery $\delta$ in normal times can be made zero, but \textbf{availability cannot}.
Labour can be reduced without limit; standing by remains. What is called ``unmanned'' is
unmanned in labour as a flow; in availability as an option it is manned.

This is a different kind of constraint from the capacity constraint of
Chapter~\ref{part:solo}. Equation~\eqref{eq:solo-capacity} bounds the \emph{total} of
delivery; availability bounds \textbf{the length of an interval over which one cannot
respond}.

\begin{definition}[Response deadline and the availability constraint]
\label{def:availability}
Write $T_{\rm resp}$ for the time allowed between the arrival of
$\omega\in\Omega\setminus\hat\Omega$ and the response, and call it the \emph{response
deadline}. Writing $u$ for the length of an interval over which the operator cannot respond,
the availability constraint is
\begin{equation}
\sup u \;\leq\; T_{\rm resp}
\label{eq:availability}
\end{equation}
\end{definition}

Equations~\eqref{eq:solo-capacity} and~\eqref{eq:availability} both have the dimension of
time, but they bind different things: the former acts on an integral, the latter on a
supremum. \textbf{Hence no reduction in the amount of labour loosens
\eqref{eq:availability}.} This offers an account of where the practical burden of a solo
business lies.

Imposing both together makes the content of that burden more concrete.

\begin{proposition}[Availability restricts the granularity of capacity]
\label{prop:granularity}
Under~\eqref{eq:availability}, the time available for delivery is fragmented into intervals
of length at most $T_{\rm resp}$. A delivery requiring an uninterrupted stretch longer than
$T_{\rm resp}$ therefore cannot be performed even when the $\mathrm{Cap}$
of~\eqref{eq:solo-capacity} has room to spare.
\end{proposition}

\begin{proof}
$\sup u\leq T_{\rm resp}$ requires that no interval without response exceed $T_{\rm resp}$.
An interval devoted to delivery is an interval without response, so its length is also at
most $T_{\rm resp}$. A delivery requiring an uninterrupted stretch longer than
$T_{\rm resp}$ does not fit into such an interval.
\end{proof}

\textbf{The availability constraint does not reduce capacity; it restricts its granularity.}
The $T_{\rm dev}$ of Chapter~\ref{ch:solo} requires long uninterrupted stretches, whereas
operation is a repetition of short responses. Proposition~\ref{prop:granularity} derives the
difficulty of combining development with operation in a solo business as the intersection of
two constraints.

\begin{remark}[They loosen in different directions]
Equation~\eqref{eq:solo-capacity} loosens by adding people.
Equation~\eqref{eq:availability} does not loosen without a relief operator, and so long as
one is alone the only route is to lengthen $T_{\rm resp}$ itself. A longer response deadline
is obtained by contractual agreement or by shrinking $\Omega\setminus\hat\Omega$.
\end{remark}

\begin{remark}[The level of $T_{\rm resp}$ differs by component]
\label{rem:tresp-components}
$T_{\rm resp}$ depends on the content of $\Omega\setminus\hat\Omega$.
Chapter~\ref{ch:platformfee} decomposes $\lambda_{\rm novel}$ into four components, and the
grace allowed for a published vulnerability differs by orders of magnitude from that for a
change in the law. Evaluating~\eqref{eq:availability} in practice requires $T_{\rm resp}$ by
component, which is not measured here.
\end{remark}

\begin{remark}[Correspondence with Knight's distinction]
\label{rem:knight}
Variation over $\hat\Omega$ has a measure defined on it and is therefore risk; both
automation and insurance (family 5-4) are possible. $\Omega\setminus\hat\Omega$ is
uncertainty, and neither is.
\end{remark}

\begin{remark}[A degree of unmannedness cannot be defined]
\label{rem:automation-rate-illposed}
``How unmanned is it?'' cannot be expressed as a ratio.

Measuring on the delivery side does not work. $\delta$ carries monetary units through the
valuation map $v$, but $v$ measures value, not effort. The ratio of the value of automated to
manual delivery does not express a degree of unmannedness.

Measuring on the response side does not work either. Automation means that $a$ is specified
in advance on $\hat\Omega$, so a degree would be the ratio of the sizes of $\hat\Omega$ and
$\Omega$. But by Remark~\ref{rem:knight} no measure is defined on
$\Omega\setminus\hat\Omega$.

\textbf{The obstacle is therefore not the absence of a unit but the absence of a measure, and
sharpening the definition does not remove it.} Unmanned operation can only be stated as a
function of duration, not of degree (Proposition~\ref{prop:unmanned-period}).
\end{remark}

\begin{remark}[Unmanned operation and profit]
If $\Phi$ and $C_{\mathrm{prod}}$ can be fully specified in advance, they are fully
replicable. Competition therefore presses $\phi_{\rm prod}\to 0$. The closer to fully
unmanned, the thinner the source of profit.

One may read profit as the return to the residual that cannot be automated. Network effects
and brand create exceptions, so this cannot be asserted flatly, but the mechanism is
plausibly at work.
\end{remark}

\subsection{Where a free period fits}

A free trial is a state in which $\delta$ is supplied with $\pi=0$, that is, the operator
extends credit to the customer, $\kappa>0$. It works because
\textbf{the marginal cost is zero, so the cost of extending that credit is near zero}. A free
trial of a physical good loses inventory; a free period for software costs only
infrastructure.

\begin{equation}
\text{free trial} = \text{credit that only a zero-marginal-cost business can issue cheaply}
\end{equation}

Section~\ref{sec:bootstrap} listed three solutions to the credit bootstrap problem, and
this is a fourth: \textbf{extend the credit oneself in place of having credit}.
$T_{\rm credit}$ is being purchased at the cost of $\kappa>0$.

\section{Decomposing the bootstrap}

\subsection{Each term and where it can be externalized}

For each term of~\eqref{eq:E-required-2}, who can absorb it. The period required for
registration with the institutional layer is added as $T_{\rm inst}$.

\begin{center}
\begin{tabularx}{\textwidth}{@{}lXXl@{}}
\toprule
Term & Content & Externalized to & Price \\
\midrule
$T_{\rm dev}$ & building the delivery system & off-the-shelf platforms, payment SDKs, hosting & usage fees \\
$T_{\rm inst}$ & registration with the institutional layer & merchant screening by a payment processor, merchant of record & fees \\
$T_{\rm credit}$ & accumulating credit & a platform's refund guarantee; or self-funded through a free period & fees, $\kappa>0$ \\
$T_{\rm acq}$ & acquiring customers & distribution through stores and marketplaces & sales commission \\
Enumerating $\hat\Omega$ & designing the envisaged states & \textbf{not possible} & --- \\
Choosing $\Phi$ & deciding the contract form & \textbf{not possible} & --- \\
\bottomrule
\end{tabularx}
\captionof{table}{The terms of the bootstrap cost, where each can be externalized, and at what price.}
\label{tab:bootstrap-outsourcing}
\end{center}

\begin{proposition}[The irreducible residue]
The only two things that cannot be externalized are the choice of $\Phi$ and the enumeration
of $\hat\Omega$.
\end{proposition}

Externalization has a consistent price: \textbf{permanently giving up $\phi_{\rm barg}$}.
A one-off launch cost is converted into a perpetual cost in the form of fees.

\subsection{The dominant term by family}

Among the types that survived in Chapter~\ref{part:solo}, the dominant component of the
launch cost differs.

\begin{center}
\begin{tabular}{@{}llllll@{}}
\toprule
Type & $T_{\rm dev}$ & $T_{\rm credit}$ & $T_{\rm acq}$ & Dominant \\
\midrule
1-1 one-off digital goods & high & low & low & $T_{\rm dev}$ \\
2-1 flat access & high & medium & medium & variance \\
6-4 charging for opportunity & medium & low & high & $T_{\rm acq}$ (two-sided) \\
7-4 donation and support & low & high & medium & $T_{\rm credit}$ \\
5-5 IP licensing & medium & high & medium & $T_{\rm credit}$ \\
\bottomrule
\end{tabular}
\captionof{table}{The dominant component of launch cost for each surviving type.}
\label{tab:startup-cost-by-type}
\end{center}

Noting that the dominant components of 1-1 and 7-4 are complementary suggests a design that
splits the launch cost as a transition across families.

\begin{center}
\begin{tabularx}{\textwidth}{@{}lX@{}}
\toprule
Stage & Content \\
\midrule
Publishing & accumulate $T_{\rm credit}$. No revenue. $T_{\rm dev}$ advances as a by-product \\
7-4 support & the first $\pi$ stands on accumulated credit alone; demands almost no $T_{\rm dev}$ \\
1-1 one-off sale & the launch is light because $T_{\rm credit}$ is done. $\kappa\approx 0$ \\
2-1 flat access & $\kappa<0$, and growth generates cash in the sense of Chapter~\ref{sec:growth} \\
\bottomrule
\end{tabularx}
\captionof{table}{Splitting the launch cost through a transition across families.}
\label{tab:phase-transition-design}
\end{center}

The leakage term $\theta$ of Chapter~\ref{part:solo} (\eqref{eq:leakage}) here works in the
opposite direction, as a contribution from publishing to $T_{\rm dev}$.
\textbf{The launch cost need not be paid all at once; it can be split by moving across
families.}

\section{Limits of this chapter}
\label{sec:unmanned-limits}

\begin{enumerate}[label=(\arabic*)]
\item \textbf{There is no quantity expressing a degree of unmannedness.}
As Remark~\ref{rem:automation-rate-illposed} states, the obstacle is the absence of a measure,
not of a unit, and sharpening the definition does not remove it. This chapter states only
what concerns duration.

\item \textbf{No level is given for $T_{\rm resp}$.}
Equation~\eqref{eq:availability} can be written, but its right-hand side differs by component
of $\Omega\setminus\hat\Omega$ and no value is identified
(Remark~\ref{rem:tresp-components}).

\item \textbf{The inverse relation between unmannedness and profit is only a mechanism noted.}
No verification, including the treatment of counterexamples such as network effects and
brand, has been carried out.

\item \textbf{No price is given for externalizing the launch cost.}
Table~\ref{tab:bootstrap-outsourcing} lists where each term can be externalized but gives no
numbers. Chapter~\ref{ch:platformfee} measures them.
\end{enumerate}
